% Generated draft: requires mathematical and source review.
\documentclass[11pt]{article}
\usepackage[margin=25mm]{geometry}
\usepackage{fontspec}
\setmainfont{Latin Modern Roman}
\usepackage{amsmath,amssymb,mathtools,unicode-math}
\setmathfont{Latin Modern Math}
\usepackage{enumitem,needspace,xurl,hyperref}
\setlength{\parindent}{0pt}
\setlength{\parskip}{5pt}
\setlength{\emergencystretch}{3em}
\newcommand{\N}{\mathbb{N}}
\newcommand{\Z}{\mathbb{Z}}
\newcommand{\Q}{\mathbb{Q}}
\newcommand{\R}{\mathbb{R}}
\newcommand{\C}{\mathbb{C}}
\newcommand{\F}{\mathbb{F}}
\newcommand{\A}{\mathbb{A}}
\newcommand{\PP}{\mathbb{P}}
\newcommand{\E}{\mathbb{E}}
\newcommand{\eps}{\varepsilon}
\newcommand{\dd}{\,\mathrm{d}}
\begin{document}
\section*{UnsolvedMath ID: 30004047}
\textbf{Source problem code:} catalogue entry
\subsection{Definitions and mathematical statement}
The supplied record leaves the parameter domains, graph conventions, and universal quantifier in the definition of \(\phi\) implicit. The following is a conditional formalization using finite simple undirected graphs, real density parameters \(x,y\in(0,1]\), positive integers \(k\), and two-edge paths whose middle vertex lies in \(B\). These conventions require confirmation against the cited source.

For such \(x,y\), let \(\mathcal G(x,y)\) consist of all tuples \((G,A,B,C)\), where \(G\) is a finite simple undirected graph and \(A,B,C\subseteq V(G)\) are pairwise disjoint nonempty sets satisfying
\[
 |N_G(a)\cap B|\ge x|B|\quad\text{for every }a\in A,
 \qquad
 |N_G(b)\cap C|\ge y|C|\quad\text{for every }b\in B.
\]
Here \(N_G(v)\) denotes the set of neighbors of \(v\). For \(c\in C\), define
\[
 R_{A,B}(c)=\{a\in A:\text{ there exists }b\in B\text{ with }ab,bc\in E(G)\}.
\]
Thus \(R_{A,B}(c)\) counts distinct vertices of \(A\), rather than paths. Interpret the universal guaranteed proportion as
\[
 \phi(x,y)=\inf_{(G,A,B,C)\in\mathcal G(x,y)}
 \max_{c\in C}\frac{|R_{A,B}(c)|}{|A|}.
\]
Equivalently, \(\phi(x,y)\) is the largest \(z\in[0,1]\) such that every tuple in \(\mathcal G(x,y)\) has some \(c\in C\) with \(|R_{A,B}(c)|\ge z|A|\). The record does not explicitly state this quantification over all admissible graphs; without it, the largest proportion would depend on the particular graph.

Under these conventions, the questions are whether
\[
 \forall x\in(1/3,1],\qquad \phi(x,x)\ge\frac12,
\]
and, more generally, whether
\[
 \forall k\in\mathbb Z_{\ge1}\ \forall x,y\in(0,1],\qquad
 \bigl(x+ky>1\ \text{and}\ kx+y\ge1\bigr)
 \Longrightarrow \phi(x,y)\ge\frac1k.
\]
The strict inequality in \(x+ky>1\) and the weak inequality in \(kx+y\ge1\) are retained exactly. The supplied record does not settle whether zero density parameters are intended to be included, or specify the intended range of \(k\).
\subsection{Short English statement}
Assume every vertex of \(A\) has neighbors in at least an \(x\)-fraction of \(B\), and every vertex of \(B\) has neighbors in at least a \(y\)-fraction of \(C\). Must some vertex of \(C\) reach at least half of \(A\) through \(B\) when \(x=y>1/3\)? More generally, do \(x+ky>1\) and \(kx+y\ge1\) guarantee reachability of at least a \(1/k\)-fraction? The displayed formalization uses explicitly provisional domain and path conventions.
\subsection{Sources}
\begin{itemize}
\item Supplied local UnsolvedMath registry record, permanent ID \(30004047\), exact source problem code \texttt{catalogue entry}, title ``Two-Step Path Density in Tripartite Graphs,'' in the Oberwolfach Reports Open Problems collection. Project registry references are \texttt{data/unsolvedmath/problems.json} and \texttt{catalog/problems.json}; the supplied record's file membership was not independently checked.
\item Supplied report citation: \emph{Graph Theory (2020)}, Oberwolfach Reports, DOI 10.4171/owr/2019/1, pp.\ 5--63. \url{https://doi.org/10.4171/owr/2019/1}. The record gives proposed year 2019 and names no proposer. It describes the collection as drawn from mathematical workshop reports and their open-problem discussions.
\item Chudnovsky et al., \emph{Concatenating bipartite graphs}, as attributed in the supplied literature sources. \url{https://doi.org/10.37236/8451}.
\item \emph{Concatenating bipartite graphs}, preprint, as listed in the supplied literature sources. \url{https://arxiv.org/abs/1902.10878}.
\item Supplied literature check date: \texttt{2026-08-22}. This is supplied dated curation, not an independently verified assessment. The record labels the problem open and reports that the published study leaves the \(x>1/3\) threshold question open while proving only a weaker lower bound. The weaker bound is not specified. The background expressly calls this a dated literature triage derived from the supplied OWR curation and requests independent verification of time-sensitive status and cited claims.
\item The supplied statement status is \texttt{corrected\_verified}; its accompanying description says mathematical content was preserved while notation, typography, duplicate clauses, or omitted definitions were repaired for self-contained publication. This supplied editorial label does not establish mathematical correctness of this draft.
\end{itemize}
\end{document}
